% ===========================================================================
% Sequencing: the layer-count minimization, its min--max theorem, and the
% public layer budget used to pad the layout.
%
% Self-contained section; drop % ===========================================================================
% Sequencing: the layer-count minimization, its min--max theorem, and the
% public layer budget used to pad the layout.
%
% Self-contained section; drop % ===========================================================================
% Sequencing: the layer-count minimization, its min--max theorem, and the
% public layer budget used to pad the layout.
%
% Self-contained section; drop % ===========================================================================
% Sequencing: the layer-count minimization, its min--max theorem, and the
% public layer budget used to pad the layout.
%
% Self-contained section; drop \input{sequencing-note} into the paper (needs
% amsmath, amssymb and a theorem environment -- llncs provides \begin{theorem}
% etc., so nothing else is required).
%
% Notation used here matches the implementation:
%   N   RLWE slot count (rows / bins)
%   b   number of blocks, m = b*N the OKVS length
%   z   span parameter (span_blocks): a layer covers z consecutive blocks
%   n   sender set size
% ===========================================================================

\section{Sequencing and the Layer Budget}
\label{sec:sequencing}

\subsection{The combinatorial problem}

After the column permutation, the $i$-th sender item occupies the
\emph{position} $\mathsf{H}_1(y_i) \in [m]$, which we read as a pair
\[
  \tau_i \;=\; \text{bin (row)} \in [N],
  \qquad
  \beta_i \;=\; \text{block} \in [b],
  \qquad m = bN .
\]
A \emph{layer} is a set of items that can be decoded by one homomorphic
multiply--accumulate chain, which imposes two constraints:

\begin{enumerate}
  \item[(C1)] \emph{one item per bin}: two items in the same layer must have
        distinct bins, since a layer contributes at most one value per slot;
  \item[(C2)] \emph{bounded span}: the blocks occupied by a layer must fit in
        $z$ consecutive blocks, since the layer's diagonals are drawn from a
        window of the encoding, and a wider window means more ciphertext
        multiplications.
\end{enumerate}

\emph{Sequencing} is the problem of partitioning the $n$ items into as few
layers as possible subject to (C1) and (C2). Write $L^\star$ for that minimum.

Both constraints are pairwise, so a layer is an independent set of the
conflict graph $G = G_{\mathrm{bin}} \cup G_{\mathrm{span}}$ ($G_{\mathrm{bin}}$
joins equal bins, $G_{\mathrm{span}}$ joins blocks farther than $z$ apart) and
$L^\star = \chi(G)$. Colouring is hopeless in general -- and indeed optimizing
over unions of two interval graphs is NP-hard -- so we use the following
reformulation instead, which is what makes the problem tractable.

\paragraph{Layers are windows.}
Constraint (C2) says exactly that the blocks of a layer lie in some window
$[s, s+z-1]$; together with (C1), a layer \emph{is} a window with capacity one
per bin. So a layering is: a multiset $W$ of window starts (one per layer),
plus an assignment of each item to a window that covers its block, injectively
within each bin. An item at block $\beta$ can use the $z$ starts
\begin{equation}
  \label{eq:starts}
  s \in [\beta - z + 1,\; \beta].
\end{equation}
Minimizing the number of layers is minimizing $|W|$.

\subsection{When does a set of windows suffice?}

Capacity is one per \emph{(window, bin)} pair, so different bins never compete:
for a fixed $W$, feasibility decomposes into one matching problem per bin. By
\eqref{eq:starts} each item's admissible starts form an interval, i.e. each
per-bin problem is a \emph{convex bipartite} matching, for which Hall's
condition collapses to intervals. Writing
\[
  \mathrm{ext}(J) \;=\; [x-z+1,\; y]
  \quad\text{for } J = [x,y],
  \qquad
  n_\tau(J) = \#\{i : \tau_i = \tau,\ \beta_i \in J\},
\]
and $W(S) = \sum_{s \in S} W(s)$, we get:

\begin{lemma}[Hall over intervals]
\label{lem:hall}
A window multiset $W$ admits all $n$ items if and only if
\[
  W\big(\mathrm{ext}(J)\big) \;\ge\; M(J) := \max_{\tau \in [N]} n_\tau(J)
  \qquad\text{for every block interval } J .
\]
\end{lemma}

\begin{proof}[Proof sketch]
Necessity is clear: the items of one bin inside $J$ can only be served by
windows starting in $\mathrm{ext}(J)$. For sufficiency, fix a bin $\tau$ and a
set $S$ of its items; split the neighbourhood $N(S)$ into maximal runs of
consecutive starts. The items of one run lie in a block interval
$J_c = [x_c, y_c]$ whose neighbourhood is exactly $\mathrm{ext}(J_c)$, and the
runs are disjoint, so the supplies add: $|N(S)| \ge \sum_c W(\mathrm{ext}(J_c))
\ge \sum_c n_\tau(J_c) \ge |S|$. Hall's theorem then gives a matching for every
bin, and the matchings for different bins do not interact.
\end{proof}

\subsection{A min--max theorem}

Lemma~\ref{lem:hall} turns sequencing into a covering problem: minimize
$\sum_s W(s)$ subject to $W(\mathrm{ext}(J)) \ge M(J)$ for all intervals $J$,
over integers $W(s) \ge 0$. Each constraint's support $\mathrm{ext}(J)$ is a
set of \emph{consecutive} starts, so the constraint matrix has the
consecutive-ones property and is therefore totally unimodular. Its LP
relaxation has an integral optimum -- so the LP value \emph{is} $L^\star$ --
and the dual is integral as well, which yields:

\begin{theorem}[Min--max for sequencing]
\label{thm:minmax}
\[
  L^\star \;=\; \max\Big\{ \textstyle\sum_{k=1}^{K} M(J_k) \;:\;
     J_1,\dots,J_K \text{ block intervals with } \mathrm{ext}(J_k)
     \text{ pairwise disjoint} \Big\}.
\]
\end{theorem}

The ``$\ge$'' direction is combinatorial and worth stating on its own, since it
is what the security argument leans on: fix such a family; a layer serving an
item of $J_k$'s heaviest bin must start in $\mathrm{ext}(J_k)$, the extensions
are disjoint, so each layer is counted for at most one $k$, and $J_k$ needs
$M(J_k)$ distinct layers. Two special cases are instructive:
\begin{itemize}
  \item $K = 1$ with $J = [1,b]$ gives $L^\star \ge \max_\tau (\text{load of }
        \tau)$, the pigeonhole bound;
  \item chopping $[1,b]$ into $\lceil b/z \rceil$ windows and adding their
        per-window maxima is one feasible \emph{primal}, i.e. an upper bound;
        it is what a naive ``hard partition'' sequencer achieves, and
        Theorem~\ref{thm:minmax} shows it is generally far from optimal.
\end{itemize}

\paragraph{The algorithm.}
We sequence with a single left-to-right first-fit pass: process blocks in
order and place each item in the admissible layer of smallest anchor, opening
a new layer anchored at the current block when none fits. It costs $O(n L)$
-- $14$\,ms at $n = 2^{20}$. (The covering program could also be solved
exactly by the standard interval-covering sweep -- add the missing windows at
the latest admissible start whenever a constraint falls short, then assign
each item to the leftmost admissible window, which is optimal for convex
bipartite matching~\cite{Glover67} -- in $O(nb + b^2)$, $130$\,ms at
$n = 2^{20}$; the point of this subsection is that nothing of the sort is
needed.) Despite its simplicity the first-fit pass is not a heuristic at all:

\begin{theorem}[The first-fit sequencer is exact]
\label{thm:greedy}
On every input the first-fit sequencer opens exactly $L^\star$ layers.
More precisely, writing $G(p)$ for the number of its layers with anchor
$\le p$ and
\[
  D(p) \;=\; \max\Big\{ \textstyle\sum_{k} M(J_k) \;:\;
     J_k = [x_k, y_k] \text{ with } \mathrm{ext}(J_k)
     \text{ pairwise disjoint},\ y_k \le p \Big\},
\]
it achieves $G(p) = D(p)$ for every $p$, whereas \emph{every} feasible
layering $W$ satisfies $\sum_{s \le p} W(s) \ge D(p)$. The first-fit layering
is therefore optimal, and pointwise the latest optimal one.
\end{theorem}

The proof rests on a structural property of the failure events. Throughout,
a layer's \emph{anchor} $a(\ell)$ is the block current when it was opened (so
anchors are created in nondecreasing order, and a layer admits blocks
$[a(\ell), a(\ell)+z-1]$), and ``$\tau$-occupied'' means the layer already
holds an item of bin $\tau$; occupied slots never vacate.

\begin{lemma}[Blocking structure]
\label{lem:blocking}
Suppose the first-fit sequencer, holding the layer multiset $\Lambda$, fails
on an item $t$ of bin $\tau$ at block $y$. Then there is an
$X \le y - z + 1$ such that
(i) the layers with $a(\ell) \ge X$ all have $a(\ell) \in [X, y]$, and each
of them is $\tau$-occupied by an item with block in $[X + z - 1,\, y]$;
(ii) hence, with $K = |\{\ell \in \Lambda : a(\ell) \ge X\}|$ and counting
$t$ itself, the items processed so far satisfy
$M([X + z - 1,\, y]) \ge K + 1$.
\end{lemma}

\begin{proof}
Build a closure. Maintain a collection of \emph{ranges} of anchors, initially
the single range $[y - z + 1,\, y]$; whenever a layer $\ell \in \Lambda$ has
its anchor inside a known range, absorb it: we argue below that $\ell$ is
$\tau$-occupied, by a unique item $w_\ell$ (its \emph{witness}) with block
$\beta_\ell$; add the range $[\beta_\ell - z + 1,\; s_\ell]$, where
$s_\ell = a(\ell)$. Iterate to a fixpoint; let $X$ be the least left endpoint
of any range.

\emph{Every absorbed layer is $\tau$-occupied.} Induct over the absorption
order. A layer absorbed by the initial range was admissible to $t$, and $t$
failed. Otherwise $\ell$'s anchor lies in $[\beta_p - z + 1, s_p]$ for an
earlier witness $p$. If $a(\ell) < s_p$: layers are opened at the block
current at their creation, so $\ell$ was opened while block
$a(\ell) < s_p \le \beta_p$ was current -- before $p$ was processed -- and at
$p$'s turn it was admissible ($a(\ell) \ge \beta_p - z + 1$) and preceded
$p$'s layer in the smallest-anchor-first scan; $p$ passing it over means it
was $\tau$-occupied then, hence still. If $a(\ell) = s_p$: either $\ell$ is
$p$'s own layer, occupied by $p$; or trace how $p$'s layer was absorbed --
through a range containing $s_p$ strictly below its endpoint, in which case
the previous case applies to $\ell$ verbatim; through the initial range, so
$a(\ell) \in [y-z+1, y]$ and $t$'s failure applies; or through another
witness's endpoint $s_q = s_p$, and we ascend to $q$. The ascent follows the
absorption order, so it terminates in one of the resolved cases.

\emph{The ranges cover exactly $[X, y]$.} A new range
$[\beta_\ell - z + 1, s_\ell]$ contains $s_\ell$ (an item's block is within
$z - 1$ of its layer's anchor), which lies in the range that absorbed
$\ell$; so the union stays a single interval, and its right end stays $y$
(no anchor exceeds the current block).

Consequently every layer with $a(\ell) \ge X$ has $a(\ell) \in [X, y]$, is
covered by a range, and was absorbed -- which is (i), since each witness has
$\beta_\ell - z + 1 \ge X$ (its range's left endpoint is one of those $X$
minimizes over), i.e. $\beta_\ell \ge X + z - 1$. For (ii): the witnesses
are distinct (one per layer, one bin-$\tau$ slot per layer), all of bin
$\tau$, with blocks in $[X + z - 1, y]$; together with $t$ (block
$y \ge X + z - 1$) they number $K + 1$.
\end{proof}

\begin{proof}[Proof of Theorem~\ref{thm:greedy}]
First, \emph{localized weak duality}: for any layering $W$ of the items
processed so far and any family as in the definition of $D(p)$, the
$M(J_k)$ items of $J_k$'s heaviest bin occupy $M(J_k)$ distinct windows with
starts in $\mathrm{ext}(J_k) \subseteq (-\infty, y_k] \subseteq
(-\infty, p]$, and the extensions are disjoint, so
$\sum_{s \le p} W(s) \ge \sum_k M(J_k)$. Hence
$\sum_{s \le p} W(s) \ge D(p)$ for every feasible $W$ -- in particular
$G(p) \ge D(p)$ at termination, since the first-fit layering is feasible.

Second, the invariant $G(p) \le D(p)$ for all $p$, by induction over the
processed items ($G$, $D$, $M$ all refer to the current prefix). Placing an
item in an existing layer changes no $G(p)$ and never decreases $D(p)$.
Suppose item $t$ (bin $\tau$, block $y$) triggers a new layer, and let $X$,
$K$ be as in Lemma~\ref{lem:blocking}. For $p < y$ nothing changes: the new
anchor is $y$. For $p \ge y$, take a family with right endpoints
$\le X - 1$ certifying $G(X-1)$ (induction), and adjoin
$J_0 = [X + z - 1,\, y]$: its extension $\mathrm{ext}(J_0) = [X, y]$ is
disjoint from the previous extensions $\subseteq (-\infty, X - 1]$, and
$M(J_0) \ge K + 1$ by the lemma, counting $t$. The family's value is now at
least $G(X - 1) + K + 1$, which by Lemma~\ref{lem:blocking}(i) is exactly
the new $G(p)$.

At termination the two directions give $G(p) = D(p)$ for every $p$; at
$p = b$, the first-fit count $G(b)$ equals $D(b)$, which by weak duality is
at most $L^\star$, while $G(b) \ge L^\star$ by feasibility.
\end{proof}

\begin{remark}
The argument re-proves the min--max theorem constructively: the first-fit
layering is feasible with $D(b)$ windows, so
Theorem~\ref{thm:minmax} follows without invoking total unimodularity.
The same equality $G \equiv D$ shows the first-fit sequencer places windows
at exactly the positions the exact covering sweep would -- they differ in
cost, not output. Both facts are corroborated mechanically: beyond the
$9{,}400$ randomized small instances checked against a brute-force chromatic
number and the production sets checked against the dual, the closure of
Lemma~\ref{lem:blocking} and the profile identity $G \equiv D$ were
machine-verified at every failure event across a further $5 \cdot 10^5$
instances, including exhaustive enumerations of all item placements for
small $(N, b, z, n)$.
\end{remark}

\subsection{The public layer budget}
\label{subsec:budget}

The realized layer count is a function of the sender's set: it is driven by
which items collide in a bin, and collisions are determined by
$\mathsf{H}_1$ evaluated on the sender's items. Transmitting $L^\star$ would
therefore leak. Instead both parties derive a \emph{public budget}
$\bar{L}$ from the parameters alone, and the sender always transmits exactly
$\bar{L}$ layers -- the $L^\star$ real ones, and $\bar{L} - L^\star$ padding
layers whose ciphertexts are a random plaintext multiple of a received
ciphertext, plus the usual mask.

Padding only goes upwards, so the sender aborts exactly when no layering fits
inside the budget, i.e. when $L^\star > \bar{L}$; Theorem~\ref{thm:budget}
bounds that event by $2^{-\lambda}$. The bound is a statement about the
\emph{minimum} layer count, so the sender has to attain it -- a sequencer
that overshot $L^\star$ would abort more often than the theorem allows. By
Theorem~\ref{thm:greedy} the first-fit sequencer attains it on every input,
so the abort event is exactly $\{L^\star > \bar{L}\}$. Conversely, when the budget is met, nothing about the realized count reaches
the wire: the receiver gets $\bar{L}$ ciphertexts and $\bar{L} \cdot N$
equality instances whatever $L^\star$ was, and which algorithm produced the
layering is immaterial.

\begin{remark}[Padding and noise]
Fixing the transmitted count is necessary but not sufficient: a real layer
accumulates $w + z$ ciphertext--plaintext products while a padding layer is a
single one, so their noise magnitudes differ by several bits, and a receiver
holding the decryption key can recover the noise and count the real layers.
Making the layers indistinguishable is the job of the noise flooding applied
before the ciphertexts are returned -- inflating the padding to match, which we
also tried, does not achieve it, since real layers already differ among
themselves by occupancy.
\end{remark}

\begin{theorem}[Layer budget]
\label{thm:budget}
Let the $n$ positions be uniform and independent over $[R]$, $R = m-w+1$, and
let
\[
  D(g) \;=\; \min\Big\{ k \;:\;
     \Pr\big[\,\mathrm{Bin}(n,\, g/R) \ge k \,\big] \;\le\;
     \frac{2^{-\lambda}}{N b^2} \Big\},
  \qquad
  \rho^\star \;=\; \max_{1 \le g \le b} \frac{D(g)+1}{g+z-1}.
\]
Then, with $\bar{L} = \lceil \rho^\star (b+z-1) \rceil$,
\[
  \Pr\big[\, L^\star > \bar{L} \,\big] \;\le\; 2^{-\lambda}.
\]
\end{theorem}

\begin{proof}
Fix a bin $\tau$, a start $x$ and a length $g$. Each item lands in bin $\tau$
with block in $[x, x+g-1]$ iff its position is one of the $g$ positions mapping
there, so $n_\tau([x,x+g-1]) \sim \mathrm{Bin}(n, \le g/R)$ and by definition of
$D(g)$ it exceeds $D(g)$ with probability at most $2^{-\lambda}/(N b^2)$. There
are at most $N b^2$ triples $(\tau, x, g)$, so by a union bound
\begin{equation}
  \label{eq:allJ}
  \Pr\big[\, \exists J:\ M(J) > D(|J|) \,\big] \;\le\; 2^{-\lambda}.
\end{equation}
Assume the complement. Lay windows down at uniform density $\rho^\star$ over
the $b+z-1$ admissible starts $s \in [-z+1, b-1]$, i.e.
$W(s) = \lfloor (s+1)\rho^\star \rfloor - \lfloor s \rho^\star \rfloor$; then
any $\ell$ consecutive starts receive at least $\rho^\star \ell - 1$ windows.
An interval $J$ with $|J| = g$ has $|\mathrm{ext}(J)| = g+z-1$, hence
\[
  W(\mathrm{ext}(J)) \;\ge\; \rho^\star (g+z-1) - 1 \;\ge\; D(g) \;\ge\; M(J),
\]
where the middle inequality is the definition of $\rho^\star$. By
Lemma~\ref{lem:hall} this $W$ is feasible, and it uses
$\lceil \rho^\star (b+z-1) \rceil = \bar{L}$ windows, so
$L^\star \le \bar{L}$.
\end{proof}

Three remarks on how the bound behaves.

\paragraph{It is computed, not extrapolated.}
$D(g)$ is an exact binomial quantile: the upper tail is summed term by term in
log space, so $\bar{L}$ is a closed-form function of $(n, N, b, z, \lambda)$
that both parties evaluate in milliseconds. Nothing here is fitted or
extrapolated from measurements, in contrast with the band-width selection of
the underlying OKVS.

\paragraph{Where the slack is.}
$\bar{L}$ overshoots the realized $L^\star$ for two reasons: the union bound
charges every one of the $Nb^2$ constraints separately, which is expensive for
short intervals where the mean count is small; and the uniform-density primal
pays the worst local ratio everywhere. Both shrink as $z$ grows -- the
denominator $g+z-1$ absorbs the union tax, and the binding constraint moves
towards $J = [1,b]$, where the union is over $N$ bins only. Measured at
$n = 2^{20}$ ($b = 269$):
\[
  \begin{array}{c|cccc}
    z & 20 & 40 & 60 & 80 \\ \hline
    \bar{L} & 399 & 296 & 266 & 255 \\
    \mathbb{E}[L^\star] & 196 & 178.5 & 175.5 & 175.5 \\
    \bar{L} / \mathbb{E}[L^\star] & 2.04 & 1.66 & 1.52 & 1.45
  \end{array}
\]
The pigeonhole floor -- the $2^{-\lambda}$ quantile of the maximum bin load,
which \emph{no} budget can go below -- is $232$ here, so at $z = 80$ the
certificate costs only $10\%$ over the information-theoretic minimum.

\paragraph{Cost of the padding.}
A padding layer costs one plaintext multiplication (versus $w + z$ for a real
layer), so the compute is driven by $L^\star$; the communication and the number
of equality instances are driven by $\bar{L}$. Increasing $z$ therefore trades
a larger per-layer decode against a smaller $\bar{L}$, and in our measurements
the second effect dominates up to $z \approx 80$.
 into the paper (needs
% amsmath, amssymb and a theorem environment -- llncs provides \begin{theorem}
% etc., so nothing else is required).
%
% Notation used here matches the implementation:
%   N   RLWE slot count (rows / bins)
%   b   number of blocks, m = b*N the OKVS length
%   z   span parameter (span_blocks): a layer covers z consecutive blocks
%   n   sender set size
% ===========================================================================

\section{Sequencing and the Layer Budget}
\label{sec:sequencing}

\subsection{The combinatorial problem}

After the column permutation, the $i$-th sender item occupies the
\emph{position} $\mathsf{H}_1(y_i) \in [m]$, which we read as a pair
\[
  \tau_i \;=\; \text{bin (row)} \in [N],
  \qquad
  \beta_i \;=\; \text{block} \in [b],
  \qquad m = bN .
\]
A \emph{layer} is a set of items that can be decoded by one homomorphic
multiply--accumulate chain, which imposes two constraints:

\begin{enumerate}
  \item[(C1)] \emph{one item per bin}: two items in the same layer must have
        distinct bins, since a layer contributes at most one value per slot;
  \item[(C2)] \emph{bounded span}: the blocks occupied by a layer must fit in
        $z$ consecutive blocks, since the layer's diagonals are drawn from a
        window of the encoding, and a wider window means more ciphertext
        multiplications.
\end{enumerate}

\emph{Sequencing} is the problem of partitioning the $n$ items into as few
layers as possible subject to (C1) and (C2). Write $L^\star$ for that minimum.

Both constraints are pairwise, so a layer is an independent set of the
conflict graph $G = G_{\mathrm{bin}} \cup G_{\mathrm{span}}$ ($G_{\mathrm{bin}}$
joins equal bins, $G_{\mathrm{span}}$ joins blocks farther than $z$ apart) and
$L^\star = \chi(G)$. Colouring is hopeless in general -- and indeed optimizing
over unions of two interval graphs is NP-hard -- so we use the following
reformulation instead, which is what makes the problem tractable.

\paragraph{Layers are windows.}
Constraint (C2) says exactly that the blocks of a layer lie in some window
$[s, s+z-1]$; together with (C1), a layer \emph{is} a window with capacity one
per bin. So a layering is: a multiset $W$ of window starts (one per layer),
plus an assignment of each item to a window that covers its block, injectively
within each bin. An item at block $\beta$ can use the $z$ starts
\begin{equation}
  \label{eq:starts}
  s \in [\beta - z + 1,\; \beta].
\end{equation}
Minimizing the number of layers is minimizing $|W|$.

\subsection{When does a set of windows suffice?}

Capacity is one per \emph{(window, bin)} pair, so different bins never compete:
for a fixed $W$, feasibility decomposes into one matching problem per bin. By
\eqref{eq:starts} each item's admissible starts form an interval, i.e. each
per-bin problem is a \emph{convex bipartite} matching, for which Hall's
condition collapses to intervals. Writing
\[
  \mathrm{ext}(J) \;=\; [x-z+1,\; y]
  \quad\text{for } J = [x,y],
  \qquad
  n_\tau(J) = \#\{i : \tau_i = \tau,\ \beta_i \in J\},
\]
and $W(S) = \sum_{s \in S} W(s)$, we get:

\begin{lemma}[Hall over intervals]
\label{lem:hall}
A window multiset $W$ admits all $n$ items if and only if
\[
  W\big(\mathrm{ext}(J)\big) \;\ge\; M(J) := \max_{\tau \in [N]} n_\tau(J)
  \qquad\text{for every block interval } J .
\]
\end{lemma}

\begin{proof}[Proof sketch]
Necessity is clear: the items of one bin inside $J$ can only be served by
windows starting in $\mathrm{ext}(J)$. For sufficiency, fix a bin $\tau$ and a
set $S$ of its items; split the neighbourhood $N(S)$ into maximal runs of
consecutive starts. The items of one run lie in a block interval
$J_c = [x_c, y_c]$ whose neighbourhood is exactly $\mathrm{ext}(J_c)$, and the
runs are disjoint, so the supplies add: $|N(S)| \ge \sum_c W(\mathrm{ext}(J_c))
\ge \sum_c n_\tau(J_c) \ge |S|$. Hall's theorem then gives a matching for every
bin, and the matchings for different bins do not interact.
\end{proof}

\subsection{A min--max theorem}

Lemma~\ref{lem:hall} turns sequencing into a covering problem: minimize
$\sum_s W(s)$ subject to $W(\mathrm{ext}(J)) \ge M(J)$ for all intervals $J$,
over integers $W(s) \ge 0$. Each constraint's support $\mathrm{ext}(J)$ is a
set of \emph{consecutive} starts, so the constraint matrix has the
consecutive-ones property and is therefore totally unimodular. Its LP
relaxation has an integral optimum -- so the LP value \emph{is} $L^\star$ --
and the dual is integral as well, which yields:

\begin{theorem}[Min--max for sequencing]
\label{thm:minmax}
\[
  L^\star \;=\; \max\Big\{ \textstyle\sum_{k=1}^{K} M(J_k) \;:\;
     J_1,\dots,J_K \text{ block intervals with } \mathrm{ext}(J_k)
     \text{ pairwise disjoint} \Big\}.
\]
\end{theorem}

The ``$\ge$'' direction is combinatorial and worth stating on its own, since it
is what the security argument leans on: fix such a family; a layer serving an
item of $J_k$'s heaviest bin must start in $\mathrm{ext}(J_k)$, the extensions
are disjoint, so each layer is counted for at most one $k$, and $J_k$ needs
$M(J_k)$ distinct layers. Two special cases are instructive:
\begin{itemize}
  \item $K = 1$ with $J = [1,b]$ gives $L^\star \ge \max_\tau (\text{load of }
        \tau)$, the pigeonhole bound;
  \item chopping $[1,b]$ into $\lceil b/z \rceil$ windows and adding their
        per-window maxima is one feasible \emph{primal}, i.e. an upper bound;
        it is what a naive ``hard partition'' sequencer achieves, and
        Theorem~\ref{thm:minmax} shows it is generally far from optimal.
\end{itemize}

\paragraph{The algorithm.}
We sequence with a single left-to-right first-fit pass: process blocks in
order and place each item in the admissible layer of smallest anchor, opening
a new layer anchored at the current block when none fits. It costs $O(n L)$
-- $14$\,ms at $n = 2^{20}$. (The covering program could also be solved
exactly by the standard interval-covering sweep -- add the missing windows at
the latest admissible start whenever a constraint falls short, then assign
each item to the leftmost admissible window, which is optimal for convex
bipartite matching~\cite{Glover67} -- in $O(nb + b^2)$, $130$\,ms at
$n = 2^{20}$; the point of this subsection is that nothing of the sort is
needed.) Despite its simplicity the first-fit pass is not a heuristic at all:

\begin{theorem}[The first-fit sequencer is exact]
\label{thm:greedy}
On every input the first-fit sequencer opens exactly $L^\star$ layers.
More precisely, writing $G(p)$ for the number of its layers with anchor
$\le p$ and
\[
  D(p) \;=\; \max\Big\{ \textstyle\sum_{k} M(J_k) \;:\;
     J_k = [x_k, y_k] \text{ with } \mathrm{ext}(J_k)
     \text{ pairwise disjoint},\ y_k \le p \Big\},
\]
it achieves $G(p) = D(p)$ for every $p$, whereas \emph{every} feasible
layering $W$ satisfies $\sum_{s \le p} W(s) \ge D(p)$. The first-fit layering
is therefore optimal, and pointwise the latest optimal one.
\end{theorem}

The proof rests on a structural property of the failure events. Throughout,
a layer's \emph{anchor} $a(\ell)$ is the block current when it was opened (so
anchors are created in nondecreasing order, and a layer admits blocks
$[a(\ell), a(\ell)+z-1]$), and ``$\tau$-occupied'' means the layer already
holds an item of bin $\tau$; occupied slots never vacate.

\begin{lemma}[Blocking structure]
\label{lem:blocking}
Suppose the first-fit sequencer, holding the layer multiset $\Lambda$, fails
on an item $t$ of bin $\tau$ at block $y$. Then there is an
$X \le y - z + 1$ such that
(i) the layers with $a(\ell) \ge X$ all have $a(\ell) \in [X, y]$, and each
of them is $\tau$-occupied by an item with block in $[X + z - 1,\, y]$;
(ii) hence, with $K = |\{\ell \in \Lambda : a(\ell) \ge X\}|$ and counting
$t$ itself, the items processed so far satisfy
$M([X + z - 1,\, y]) \ge K + 1$.
\end{lemma}

\begin{proof}
Build a closure. Maintain a collection of \emph{ranges} of anchors, initially
the single range $[y - z + 1,\, y]$; whenever a layer $\ell \in \Lambda$ has
its anchor inside a known range, absorb it: we argue below that $\ell$ is
$\tau$-occupied, by a unique item $w_\ell$ (its \emph{witness}) with block
$\beta_\ell$; add the range $[\beta_\ell - z + 1,\; s_\ell]$, where
$s_\ell = a(\ell)$. Iterate to a fixpoint; let $X$ be the least left endpoint
of any range.

\emph{Every absorbed layer is $\tau$-occupied.} Induct over the absorption
order. A layer absorbed by the initial range was admissible to $t$, and $t$
failed. Otherwise $\ell$'s anchor lies in $[\beta_p - z + 1, s_p]$ for an
earlier witness $p$. If $a(\ell) < s_p$: layers are opened at the block
current at their creation, so $\ell$ was opened while block
$a(\ell) < s_p \le \beta_p$ was current -- before $p$ was processed -- and at
$p$'s turn it was admissible ($a(\ell) \ge \beta_p - z + 1$) and preceded
$p$'s layer in the smallest-anchor-first scan; $p$ passing it over means it
was $\tau$-occupied then, hence still. If $a(\ell) = s_p$: either $\ell$ is
$p$'s own layer, occupied by $p$; or trace how $p$'s layer was absorbed --
through a range containing $s_p$ strictly below its endpoint, in which case
the previous case applies to $\ell$ verbatim; through the initial range, so
$a(\ell) \in [y-z+1, y]$ and $t$'s failure applies; or through another
witness's endpoint $s_q = s_p$, and we ascend to $q$. The ascent follows the
absorption order, so it terminates in one of the resolved cases.

\emph{The ranges cover exactly $[X, y]$.} A new range
$[\beta_\ell - z + 1, s_\ell]$ contains $s_\ell$ (an item's block is within
$z - 1$ of its layer's anchor), which lies in the range that absorbed
$\ell$; so the union stays a single interval, and its right end stays $y$
(no anchor exceeds the current block).

Consequently every layer with $a(\ell) \ge X$ has $a(\ell) \in [X, y]$, is
covered by a range, and was absorbed -- which is (i), since each witness has
$\beta_\ell - z + 1 \ge X$ (its range's left endpoint is one of those $X$
minimizes over), i.e. $\beta_\ell \ge X + z - 1$. For (ii): the witnesses
are distinct (one per layer, one bin-$\tau$ slot per layer), all of bin
$\tau$, with blocks in $[X + z - 1, y]$; together with $t$ (block
$y \ge X + z - 1$) they number $K + 1$.
\end{proof}

\begin{proof}[Proof of Theorem~\ref{thm:greedy}]
First, \emph{localized weak duality}: for any layering $W$ of the items
processed so far and any family as in the definition of $D(p)$, the
$M(J_k)$ items of $J_k$'s heaviest bin occupy $M(J_k)$ distinct windows with
starts in $\mathrm{ext}(J_k) \subseteq (-\infty, y_k] \subseteq
(-\infty, p]$, and the extensions are disjoint, so
$\sum_{s \le p} W(s) \ge \sum_k M(J_k)$. Hence
$\sum_{s \le p} W(s) \ge D(p)$ for every feasible $W$ -- in particular
$G(p) \ge D(p)$ at termination, since the first-fit layering is feasible.

Second, the invariant $G(p) \le D(p)$ for all $p$, by induction over the
processed items ($G$, $D$, $M$ all refer to the current prefix). Placing an
item in an existing layer changes no $G(p)$ and never decreases $D(p)$.
Suppose item $t$ (bin $\tau$, block $y$) triggers a new layer, and let $X$,
$K$ be as in Lemma~\ref{lem:blocking}. For $p < y$ nothing changes: the new
anchor is $y$. For $p \ge y$, take a family with right endpoints
$\le X - 1$ certifying $G(X-1)$ (induction), and adjoin
$J_0 = [X + z - 1,\, y]$: its extension $\mathrm{ext}(J_0) = [X, y]$ is
disjoint from the previous extensions $\subseteq (-\infty, X - 1]$, and
$M(J_0) \ge K + 1$ by the lemma, counting $t$. The family's value is now at
least $G(X - 1) + K + 1$, which by Lemma~\ref{lem:blocking}(i) is exactly
the new $G(p)$.

At termination the two directions give $G(p) = D(p)$ for every $p$; at
$p = b$, the first-fit count $G(b)$ equals $D(b)$, which by weak duality is
at most $L^\star$, while $G(b) \ge L^\star$ by feasibility.
\end{proof}

\begin{remark}
The argument re-proves the min--max theorem constructively: the first-fit
layering is feasible with $D(b)$ windows, so
Theorem~\ref{thm:minmax} follows without invoking total unimodularity.
The same equality $G \equiv D$ shows the first-fit sequencer places windows
at exactly the positions the exact covering sweep would -- they differ in
cost, not output. Both facts are corroborated mechanically: beyond the
$9{,}400$ randomized small instances checked against a brute-force chromatic
number and the production sets checked against the dual, the closure of
Lemma~\ref{lem:blocking} and the profile identity $G \equiv D$ were
machine-verified at every failure event across a further $5 \cdot 10^5$
instances, including exhaustive enumerations of all item placements for
small $(N, b, z, n)$.
\end{remark}

\subsection{The public layer budget}
\label{subsec:budget}

The realized layer count is a function of the sender's set: it is driven by
which items collide in a bin, and collisions are determined by
$\mathsf{H}_1$ evaluated on the sender's items. Transmitting $L^\star$ would
therefore leak. Instead both parties derive a \emph{public budget}
$\bar{L}$ from the parameters alone, and the sender always transmits exactly
$\bar{L}$ layers -- the $L^\star$ real ones, and $\bar{L} - L^\star$ padding
layers whose ciphertexts are a random plaintext multiple of a received
ciphertext, plus the usual mask.

Padding only goes upwards, so the sender aborts exactly when no layering fits
inside the budget, i.e. when $L^\star > \bar{L}$; Theorem~\ref{thm:budget}
bounds that event by $2^{-\lambda}$. The bound is a statement about the
\emph{minimum} layer count, so the sender has to attain it -- a sequencer
that overshot $L^\star$ would abort more often than the theorem allows. By
Theorem~\ref{thm:greedy} the first-fit sequencer attains it on every input,
so the abort event is exactly $\{L^\star > \bar{L}\}$. Conversely, when the budget is met, nothing about the realized count reaches
the wire: the receiver gets $\bar{L}$ ciphertexts and $\bar{L} \cdot N$
equality instances whatever $L^\star$ was, and which algorithm produced the
layering is immaterial.

\begin{remark}[Padding and noise]
Fixing the transmitted count is necessary but not sufficient: a real layer
accumulates $w + z$ ciphertext--plaintext products while a padding layer is a
single one, so their noise magnitudes differ by several bits, and a receiver
holding the decryption key can recover the noise and count the real layers.
Making the layers indistinguishable is the job of the noise flooding applied
before the ciphertexts are returned -- inflating the padding to match, which we
also tried, does not achieve it, since real layers already differ among
themselves by occupancy.
\end{remark}

\begin{theorem}[Layer budget]
\label{thm:budget}
Let the $n$ positions be uniform and independent over $[R]$, $R = m-w+1$, and
let
\[
  D(g) \;=\; \min\Big\{ k \;:\;
     \Pr\big[\,\mathrm{Bin}(n,\, g/R) \ge k \,\big] \;\le\;
     \frac{2^{-\lambda}}{N b^2} \Big\},
  \qquad
  \rho^\star \;=\; \max_{1 \le g \le b} \frac{D(g)+1}{g+z-1}.
\]
Then, with $\bar{L} = \lceil \rho^\star (b+z-1) \rceil$,
\[
  \Pr\big[\, L^\star > \bar{L} \,\big] \;\le\; 2^{-\lambda}.
\]
\end{theorem}

\begin{proof}
Fix a bin $\tau$, a start $x$ and a length $g$. Each item lands in bin $\tau$
with block in $[x, x+g-1]$ iff its position is one of the $g$ positions mapping
there, so $n_\tau([x,x+g-1]) \sim \mathrm{Bin}(n, \le g/R)$ and by definition of
$D(g)$ it exceeds $D(g)$ with probability at most $2^{-\lambda}/(N b^2)$. There
are at most $N b^2$ triples $(\tau, x, g)$, so by a union bound
\begin{equation}
  \label{eq:allJ}
  \Pr\big[\, \exists J:\ M(J) > D(|J|) \,\big] \;\le\; 2^{-\lambda}.
\end{equation}
Assume the complement. Lay windows down at uniform density $\rho^\star$ over
the $b+z-1$ admissible starts $s \in [-z+1, b-1]$, i.e.
$W(s) = \lfloor (s+1)\rho^\star \rfloor - \lfloor s \rho^\star \rfloor$; then
any $\ell$ consecutive starts receive at least $\rho^\star \ell - 1$ windows.
An interval $J$ with $|J| = g$ has $|\mathrm{ext}(J)| = g+z-1$, hence
\[
  W(\mathrm{ext}(J)) \;\ge\; \rho^\star (g+z-1) - 1 \;\ge\; D(g) \;\ge\; M(J),
\]
where the middle inequality is the definition of $\rho^\star$. By
Lemma~\ref{lem:hall} this $W$ is feasible, and it uses
$\lceil \rho^\star (b+z-1) \rceil = \bar{L}$ windows, so
$L^\star \le \bar{L}$.
\end{proof}

Three remarks on how the bound behaves.

\paragraph{It is computed, not extrapolated.}
$D(g)$ is an exact binomial quantile: the upper tail is summed term by term in
log space, so $\bar{L}$ is a closed-form function of $(n, N, b, z, \lambda)$
that both parties evaluate in milliseconds. Nothing here is fitted or
extrapolated from measurements, in contrast with the band-width selection of
the underlying OKVS.

\paragraph{Where the slack is.}
$\bar{L}$ overshoots the realized $L^\star$ for two reasons: the union bound
charges every one of the $Nb^2$ constraints separately, which is expensive for
short intervals where the mean count is small; and the uniform-density primal
pays the worst local ratio everywhere. Both shrink as $z$ grows -- the
denominator $g+z-1$ absorbs the union tax, and the binding constraint moves
towards $J = [1,b]$, where the union is over $N$ bins only. Measured at
$n = 2^{20}$ ($b = 269$):
\[
  \begin{array}{c|cccc}
    z & 20 & 40 & 60 & 80 \\ \hline
    \bar{L} & 399 & 296 & 266 & 255 \\
    \mathbb{E}[L^\star] & 196 & 178.5 & 175.5 & 175.5 \\
    \bar{L} / \mathbb{E}[L^\star] & 2.04 & 1.66 & 1.52 & 1.45
  \end{array}
\]
The pigeonhole floor -- the $2^{-\lambda}$ quantile of the maximum bin load,
which \emph{no} budget can go below -- is $232$ here, so at $z = 80$ the
certificate costs only $10\%$ over the information-theoretic minimum.

\paragraph{Cost of the padding.}
A padding layer costs one plaintext multiplication (versus $w + z$ for a real
layer), so the compute is driven by $L^\star$; the communication and the number
of equality instances are driven by $\bar{L}$. Increasing $z$ therefore trades
a larger per-layer decode against a smaller $\bar{L}$, and in our measurements
the second effect dominates up to $z \approx 80$.
 into the paper (needs
% amsmath, amssymb and a theorem environment -- llncs provides \begin{theorem}
% etc., so nothing else is required).
%
% Notation used here matches the implementation:
%   N   RLWE slot count (rows / bins)
%   b   number of blocks, m = b*N the OKVS length
%   z   span parameter (span_blocks): a layer covers z consecutive blocks
%   n   sender set size
% ===========================================================================

\section{Sequencing and the Layer Budget}
\label{sec:sequencing}

\subsection{The combinatorial problem}

After the column permutation, the $i$-th sender item occupies the
\emph{position} $\mathsf{H}_1(y_i) \in [m]$, which we read as a pair
\[
  \tau_i \;=\; \text{bin (row)} \in [N],
  \qquad
  \beta_i \;=\; \text{block} \in [b],
  \qquad m = bN .
\]
A \emph{layer} is a set of items that can be decoded by one homomorphic
multiply--accumulate chain, which imposes two constraints:

\begin{enumerate}
  \item[(C1)] \emph{one item per bin}: two items in the same layer must have
        distinct bins, since a layer contributes at most one value per slot;
  \item[(C2)] \emph{bounded span}: the blocks occupied by a layer must fit in
        $z$ consecutive blocks, since the layer's diagonals are drawn from a
        window of the encoding, and a wider window means more ciphertext
        multiplications.
\end{enumerate}

\emph{Sequencing} is the problem of partitioning the $n$ items into as few
layers as possible subject to (C1) and (C2). Write $L^\star$ for that minimum.

Both constraints are pairwise, so a layer is an independent set of the
conflict graph $G = G_{\mathrm{bin}} \cup G_{\mathrm{span}}$ ($G_{\mathrm{bin}}$
joins equal bins, $G_{\mathrm{span}}$ joins blocks farther than $z$ apart) and
$L^\star = \chi(G)$. Colouring is hopeless in general -- and indeed optimizing
over unions of two interval graphs is NP-hard -- so we use the following
reformulation instead, which is what makes the problem tractable.

\paragraph{Layers are windows.}
Constraint (C2) says exactly that the blocks of a layer lie in some window
$[s, s+z-1]$; together with (C1), a layer \emph{is} a window with capacity one
per bin. So a layering is: a multiset $W$ of window starts (one per layer),
plus an assignment of each item to a window that covers its block, injectively
within each bin. An item at block $\beta$ can use the $z$ starts
\begin{equation}
  \label{eq:starts}
  s \in [\beta - z + 1,\; \beta].
\end{equation}
Minimizing the number of layers is minimizing $|W|$.

\subsection{When does a set of windows suffice?}

Capacity is one per \emph{(window, bin)} pair, so different bins never compete:
for a fixed $W$, feasibility decomposes into one matching problem per bin. By
\eqref{eq:starts} each item's admissible starts form an interval, i.e. each
per-bin problem is a \emph{convex bipartite} matching, for which Hall's
condition collapses to intervals. Writing
\[
  \mathrm{ext}(J) \;=\; [x-z+1,\; y]
  \quad\text{for } J = [x,y],
  \qquad
  n_\tau(J) = \#\{i : \tau_i = \tau,\ \beta_i \in J\},
\]
and $W(S) = \sum_{s \in S} W(s)$, we get:

\begin{lemma}[Hall over intervals]
\label{lem:hall}
A window multiset $W$ admits all $n$ items if and only if
\[
  W\big(\mathrm{ext}(J)\big) \;\ge\; M(J) := \max_{\tau \in [N]} n_\tau(J)
  \qquad\text{for every block interval } J .
\]
\end{lemma}

\begin{proof}[Proof sketch]
Necessity is clear: the items of one bin inside $J$ can only be served by
windows starting in $\mathrm{ext}(J)$. For sufficiency, fix a bin $\tau$ and a
set $S$ of its items; split the neighbourhood $N(S)$ into maximal runs of
consecutive starts. The items of one run lie in a block interval
$J_c = [x_c, y_c]$ whose neighbourhood is exactly $\mathrm{ext}(J_c)$, and the
runs are disjoint, so the supplies add: $|N(S)| \ge \sum_c W(\mathrm{ext}(J_c))
\ge \sum_c n_\tau(J_c) \ge |S|$. Hall's theorem then gives a matching for every
bin, and the matchings for different bins do not interact.
\end{proof}

\subsection{A min--max theorem}

Lemma~\ref{lem:hall} turns sequencing into a covering problem: minimize
$\sum_s W(s)$ subject to $W(\mathrm{ext}(J)) \ge M(J)$ for all intervals $J$,
over integers $W(s) \ge 0$. Each constraint's support $\mathrm{ext}(J)$ is a
set of \emph{consecutive} starts, so the constraint matrix has the
consecutive-ones property and is therefore totally unimodular. Its LP
relaxation has an integral optimum -- so the LP value \emph{is} $L^\star$ --
and the dual is integral as well, which yields:

\begin{theorem}[Min--max for sequencing]
\label{thm:minmax}
\[
  L^\star \;=\; \max\Big\{ \textstyle\sum_{k=1}^{K} M(J_k) \;:\;
     J_1,\dots,J_K \text{ block intervals with } \mathrm{ext}(J_k)
     \text{ pairwise disjoint} \Big\}.
\]
\end{theorem}

The ``$\ge$'' direction is combinatorial and worth stating on its own, since it
is what the security argument leans on: fix such a family; a layer serving an
item of $J_k$'s heaviest bin must start in $\mathrm{ext}(J_k)$, the extensions
are disjoint, so each layer is counted for at most one $k$, and $J_k$ needs
$M(J_k)$ distinct layers. Two special cases are instructive:
\begin{itemize}
  \item $K = 1$ with $J = [1,b]$ gives $L^\star \ge \max_\tau (\text{load of }
        \tau)$, the pigeonhole bound;
  \item chopping $[1,b]$ into $\lceil b/z \rceil$ windows and adding their
        per-window maxima is one feasible \emph{primal}, i.e. an upper bound;
        it is what a naive ``hard partition'' sequencer achieves, and
        Theorem~\ref{thm:minmax} shows it is generally far from optimal.
\end{itemize}

\paragraph{The algorithm.}
We sequence with a single left-to-right first-fit pass: process blocks in
order and place each item in the admissible layer of smallest anchor, opening
a new layer anchored at the current block when none fits. It costs $O(n L)$
-- $14$\,ms at $n = 2^{20}$. (The covering program could also be solved
exactly by the standard interval-covering sweep -- add the missing windows at
the latest admissible start whenever a constraint falls short, then assign
each item to the leftmost admissible window, which is optimal for convex
bipartite matching~\cite{Glover67} -- in $O(nb + b^2)$, $130$\,ms at
$n = 2^{20}$; the point of this subsection is that nothing of the sort is
needed.) Despite its simplicity the first-fit pass is not a heuristic at all:

\begin{theorem}[The first-fit sequencer is exact]
\label{thm:greedy}
On every input the first-fit sequencer opens exactly $L^\star$ layers.
More precisely, writing $G(p)$ for the number of its layers with anchor
$\le p$ and
\[
  D(p) \;=\; \max\Big\{ \textstyle\sum_{k} M(J_k) \;:\;
     J_k = [x_k, y_k] \text{ with } \mathrm{ext}(J_k)
     \text{ pairwise disjoint},\ y_k \le p \Big\},
\]
it achieves $G(p) = D(p)$ for every $p$, whereas \emph{every} feasible
layering $W$ satisfies $\sum_{s \le p} W(s) \ge D(p)$. The first-fit layering
is therefore optimal, and pointwise the latest optimal one.
\end{theorem}

The proof rests on a structural property of the failure events. Throughout,
a layer's \emph{anchor} $a(\ell)$ is the block current when it was opened (so
anchors are created in nondecreasing order, and a layer admits blocks
$[a(\ell), a(\ell)+z-1]$), and ``$\tau$-occupied'' means the layer already
holds an item of bin $\tau$; occupied slots never vacate.

\begin{lemma}[Blocking structure]
\label{lem:blocking}
Suppose the first-fit sequencer, holding the layer multiset $\Lambda$, fails
on an item $t$ of bin $\tau$ at block $y$. Then there is an
$X \le y - z + 1$ such that
(i) the layers with $a(\ell) \ge X$ all have $a(\ell) \in [X, y]$, and each
of them is $\tau$-occupied by an item with block in $[X + z - 1,\, y]$;
(ii) hence, with $K = |\{\ell \in \Lambda : a(\ell) \ge X\}|$ and counting
$t$ itself, the items processed so far satisfy
$M([X + z - 1,\, y]) \ge K + 1$.
\end{lemma}

\begin{proof}
Build a closure. Maintain a collection of \emph{ranges} of anchors, initially
the single range $[y - z + 1,\, y]$; whenever a layer $\ell \in \Lambda$ has
its anchor inside a known range, absorb it: we argue below that $\ell$ is
$\tau$-occupied, by a unique item $w_\ell$ (its \emph{witness}) with block
$\beta_\ell$; add the range $[\beta_\ell - z + 1,\; s_\ell]$, where
$s_\ell = a(\ell)$. Iterate to a fixpoint; let $X$ be the least left endpoint
of any range.

\emph{Every absorbed layer is $\tau$-occupied.} Induct over the absorption
order. A layer absorbed by the initial range was admissible to $t$, and $t$
failed. Otherwise $\ell$'s anchor lies in $[\beta_p - z + 1, s_p]$ for an
earlier witness $p$. If $a(\ell) < s_p$: layers are opened at the block
current at their creation, so $\ell$ was opened while block
$a(\ell) < s_p \le \beta_p$ was current -- before $p$ was processed -- and at
$p$'s turn it was admissible ($a(\ell) \ge \beta_p - z + 1$) and preceded
$p$'s layer in the smallest-anchor-first scan; $p$ passing it over means it
was $\tau$-occupied then, hence still. If $a(\ell) = s_p$: either $\ell$ is
$p$'s own layer, occupied by $p$; or trace how $p$'s layer was absorbed --
through a range containing $s_p$ strictly below its endpoint, in which case
the previous case applies to $\ell$ verbatim; through the initial range, so
$a(\ell) \in [y-z+1, y]$ and $t$'s failure applies; or through another
witness's endpoint $s_q = s_p$, and we ascend to $q$. The ascent follows the
absorption order, so it terminates in one of the resolved cases.

\emph{The ranges cover exactly $[X, y]$.} A new range
$[\beta_\ell - z + 1, s_\ell]$ contains $s_\ell$ (an item's block is within
$z - 1$ of its layer's anchor), which lies in the range that absorbed
$\ell$; so the union stays a single interval, and its right end stays $y$
(no anchor exceeds the current block).

Consequently every layer with $a(\ell) \ge X$ has $a(\ell) \in [X, y]$, is
covered by a range, and was absorbed -- which is (i), since each witness has
$\beta_\ell - z + 1 \ge X$ (its range's left endpoint is one of those $X$
minimizes over), i.e. $\beta_\ell \ge X + z - 1$. For (ii): the witnesses
are distinct (one per layer, one bin-$\tau$ slot per layer), all of bin
$\tau$, with blocks in $[X + z - 1, y]$; together with $t$ (block
$y \ge X + z - 1$) they number $K + 1$.
\end{proof}

\begin{proof}[Proof of Theorem~\ref{thm:greedy}]
First, \emph{localized weak duality}: for any layering $W$ of the items
processed so far and any family as in the definition of $D(p)$, the
$M(J_k)$ items of $J_k$'s heaviest bin occupy $M(J_k)$ distinct windows with
starts in $\mathrm{ext}(J_k) \subseteq (-\infty, y_k] \subseteq
(-\infty, p]$, and the extensions are disjoint, so
$\sum_{s \le p} W(s) \ge \sum_k M(J_k)$. Hence
$\sum_{s \le p} W(s) \ge D(p)$ for every feasible $W$ -- in particular
$G(p) \ge D(p)$ at termination, since the first-fit layering is feasible.

Second, the invariant $G(p) \le D(p)$ for all $p$, by induction over the
processed items ($G$, $D$, $M$ all refer to the current prefix). Placing an
item in an existing layer changes no $G(p)$ and never decreases $D(p)$.
Suppose item $t$ (bin $\tau$, block $y$) triggers a new layer, and let $X$,
$K$ be as in Lemma~\ref{lem:blocking}. For $p < y$ nothing changes: the new
anchor is $y$. For $p \ge y$, take a family with right endpoints
$\le X - 1$ certifying $G(X-1)$ (induction), and adjoin
$J_0 = [X + z - 1,\, y]$: its extension $\mathrm{ext}(J_0) = [X, y]$ is
disjoint from the previous extensions $\subseteq (-\infty, X - 1]$, and
$M(J_0) \ge K + 1$ by the lemma, counting $t$. The family's value is now at
least $G(X - 1) + K + 1$, which by Lemma~\ref{lem:blocking}(i) is exactly
the new $G(p)$.

At termination the two directions give $G(p) = D(p)$ for every $p$; at
$p = b$, the first-fit count $G(b)$ equals $D(b)$, which by weak duality is
at most $L^\star$, while $G(b) \ge L^\star$ by feasibility.
\end{proof}

\begin{remark}
The argument re-proves the min--max theorem constructively: the first-fit
layering is feasible with $D(b)$ windows, so
Theorem~\ref{thm:minmax} follows without invoking total unimodularity.
The same equality $G \equiv D$ shows the first-fit sequencer places windows
at exactly the positions the exact covering sweep would -- they differ in
cost, not output. Both facts are corroborated mechanically: beyond the
$9{,}400$ randomized small instances checked against a brute-force chromatic
number and the production sets checked against the dual, the closure of
Lemma~\ref{lem:blocking} and the profile identity $G \equiv D$ were
machine-verified at every failure event across a further $5 \cdot 10^5$
instances, including exhaustive enumerations of all item placements for
small $(N, b, z, n)$.
\end{remark}

\subsection{The public layer budget}
\label{subsec:budget}

The realized layer count is a function of the sender's set: it is driven by
which items collide in a bin, and collisions are determined by
$\mathsf{H}_1$ evaluated on the sender's items. Transmitting $L^\star$ would
therefore leak. Instead both parties derive a \emph{public budget}
$\bar{L}$ from the parameters alone, and the sender always transmits exactly
$\bar{L}$ layers -- the $L^\star$ real ones, and $\bar{L} - L^\star$ padding
layers whose ciphertexts are a random plaintext multiple of a received
ciphertext, plus the usual mask.

Padding only goes upwards, so the sender aborts exactly when no layering fits
inside the budget, i.e. when $L^\star > \bar{L}$; Theorem~\ref{thm:budget}
bounds that event by $2^{-\lambda}$. The bound is a statement about the
\emph{minimum} layer count, so the sender has to attain it -- a sequencer
that overshot $L^\star$ would abort more often than the theorem allows. By
Theorem~\ref{thm:greedy} the first-fit sequencer attains it on every input,
so the abort event is exactly $\{L^\star > \bar{L}\}$. Conversely, when the budget is met, nothing about the realized count reaches
the wire: the receiver gets $\bar{L}$ ciphertexts and $\bar{L} \cdot N$
equality instances whatever $L^\star$ was, and which algorithm produced the
layering is immaterial.

\begin{remark}[Padding and noise]
Fixing the transmitted count is necessary but not sufficient: a real layer
accumulates $w + z$ ciphertext--plaintext products while a padding layer is a
single one, so their noise magnitudes differ by several bits, and a receiver
holding the decryption key can recover the noise and count the real layers.
Making the layers indistinguishable is the job of the noise flooding applied
before the ciphertexts are returned -- inflating the padding to match, which we
also tried, does not achieve it, since real layers already differ among
themselves by occupancy.
\end{remark}

\begin{theorem}[Layer budget]
\label{thm:budget}
Let the $n$ positions be uniform and independent over $[R]$, $R = m-w+1$, and
let
\[
  D(g) \;=\; \min\Big\{ k \;:\;
     \Pr\big[\,\mathrm{Bin}(n,\, g/R) \ge k \,\big] \;\le\;
     \frac{2^{-\lambda}}{N b^2} \Big\},
  \qquad
  \rho^\star \;=\; \max_{1 \le g \le b} \frac{D(g)+1}{g+z-1}.
\]
Then, with $\bar{L} = \lceil \rho^\star (b+z-1) \rceil$,
\[
  \Pr\big[\, L^\star > \bar{L} \,\big] \;\le\; 2^{-\lambda}.
\]
\end{theorem}

\begin{proof}
Fix a bin $\tau$, a start $x$ and a length $g$. Each item lands in bin $\tau$
with block in $[x, x+g-1]$ iff its position is one of the $g$ positions mapping
there, so $n_\tau([x,x+g-1]) \sim \mathrm{Bin}(n, \le g/R)$ and by definition of
$D(g)$ it exceeds $D(g)$ with probability at most $2^{-\lambda}/(N b^2)$. There
are at most $N b^2$ triples $(\tau, x, g)$, so by a union bound
\begin{equation}
  \label{eq:allJ}
  \Pr\big[\, \exists J:\ M(J) > D(|J|) \,\big] \;\le\; 2^{-\lambda}.
\end{equation}
Assume the complement. Lay windows down at uniform density $\rho^\star$ over
the $b+z-1$ admissible starts $s \in [-z+1, b-1]$, i.e.
$W(s) = \lfloor (s+1)\rho^\star \rfloor - \lfloor s \rho^\star \rfloor$; then
any $\ell$ consecutive starts receive at least $\rho^\star \ell - 1$ windows.
An interval $J$ with $|J| = g$ has $|\mathrm{ext}(J)| = g+z-1$, hence
\[
  W(\mathrm{ext}(J)) \;\ge\; \rho^\star (g+z-1) - 1 \;\ge\; D(g) \;\ge\; M(J),
\]
where the middle inequality is the definition of $\rho^\star$. By
Lemma~\ref{lem:hall} this $W$ is feasible, and it uses
$\lceil \rho^\star (b+z-1) \rceil = \bar{L}$ windows, so
$L^\star \le \bar{L}$.
\end{proof}

Three remarks on how the bound behaves.

\paragraph{It is computed, not extrapolated.}
$D(g)$ is an exact binomial quantile: the upper tail is summed term by term in
log space, so $\bar{L}$ is a closed-form function of $(n, N, b, z, \lambda)$
that both parties evaluate in milliseconds. Nothing here is fitted or
extrapolated from measurements, in contrast with the band-width selection of
the underlying OKVS.

\paragraph{Where the slack is.}
$\bar{L}$ overshoots the realized $L^\star$ for two reasons: the union bound
charges every one of the $Nb^2$ constraints separately, which is expensive for
short intervals where the mean count is small; and the uniform-density primal
pays the worst local ratio everywhere. Both shrink as $z$ grows -- the
denominator $g+z-1$ absorbs the union tax, and the binding constraint moves
towards $J = [1,b]$, where the union is over $N$ bins only. Measured at
$n = 2^{20}$ ($b = 269$):
\[
  \begin{array}{c|cccc}
    z & 20 & 40 & 60 & 80 \\ \hline
    \bar{L} & 399 & 296 & 266 & 255 \\
    \mathbb{E}[L^\star] & 196 & 178.5 & 175.5 & 175.5 \\
    \bar{L} / \mathbb{E}[L^\star] & 2.04 & 1.66 & 1.52 & 1.45
  \end{array}
\]
The pigeonhole floor -- the $2^{-\lambda}$ quantile of the maximum bin load,
which \emph{no} budget can go below -- is $232$ here, so at $z = 80$ the
certificate costs only $10\%$ over the information-theoretic minimum.

\paragraph{Cost of the padding.}
A padding layer costs one plaintext multiplication (versus $w + z$ for a real
layer), so the compute is driven by $L^\star$; the communication and the number
of equality instances are driven by $\bar{L}$. Increasing $z$ therefore trades
a larger per-layer decode against a smaller $\bar{L}$, and in our measurements
the second effect dominates up to $z \approx 80$.
 into the paper (needs
% amsmath, amssymb and a theorem environment -- llncs provides \begin{theorem}
% etc., so nothing else is required).
%
% Notation used here matches the implementation:
%   N   RLWE slot count (rows / bins)
%   b   number of blocks, m = b*N the OKVS length
%   z   span parameter (span_blocks): a layer covers z consecutive blocks
%   n   sender set size
% ===========================================================================

\section{Sequencing and the Layer Budget}
\label{sec:sequencing}

\subsection{The combinatorial problem}

After the column permutation, the $i$-th sender item occupies the
\emph{position} $\mathsf{H}_1(y_i) \in [m]$, which we read as a pair
\[
  \tau_i \;=\; \text{bin (row)} \in [N],
  \qquad
  \beta_i \;=\; \text{block} \in [b],
  \qquad m = bN .
\]
A \emph{layer} is a set of items that can be decoded by one homomorphic
multiply--accumulate chain, which imposes two constraints:

\begin{enumerate}
  \item[(C1)] \emph{one item per bin}: two items in the same layer must have
        distinct bins, since a layer contributes at most one value per slot;
  \item[(C2)] \emph{bounded span}: the blocks occupied by a layer must fit in
        $z$ consecutive blocks, since the layer's diagonals are drawn from a
        window of the encoding, and a wider window means more ciphertext
        multiplications.
\end{enumerate}

\emph{Sequencing} is the problem of partitioning the $n$ items into as few
layers as possible subject to (C1) and (C2). Write $L^\star$ for that minimum.

Both constraints are pairwise, so a layer is an independent set of the
conflict graph $G = G_{\mathrm{bin}} \cup G_{\mathrm{span}}$ ($G_{\mathrm{bin}}$
joins equal bins, $G_{\mathrm{span}}$ joins blocks farther than $z$ apart) and
$L^\star = \chi(G)$. Colouring is hopeless in general -- and indeed optimizing
over unions of two interval graphs is NP-hard -- so we use the following
reformulation instead, which is what makes the problem tractable.

\paragraph{Layers are windows.}
Constraint (C2) says exactly that the blocks of a layer lie in some window
$[s, s+z-1]$; together with (C1), a layer \emph{is} a window with capacity one
per bin. So a layering is: a multiset $W$ of window starts (one per layer),
plus an assignment of each item to a window that covers its block, injectively
within each bin. An item at block $\beta$ can use the $z$ starts
\begin{equation}
  \label{eq:starts}
  s \in [\beta - z + 1,\; \beta].
\end{equation}
Minimizing the number of layers is minimizing $|W|$.

\subsection{When does a set of windows suffice?}

Capacity is one per \emph{(window, bin)} pair, so different bins never compete:
for a fixed $W$, feasibility decomposes into one matching problem per bin. By
\eqref{eq:starts} each item's admissible starts form an interval, i.e. each
per-bin problem is a \emph{convex bipartite} matching, for which Hall's
condition collapses to intervals. Writing
\[
  \mathrm{ext}(J) \;=\; [x-z+1,\; y]
  \quad\text{for } J = [x,y],
  \qquad
  n_\tau(J) = \#\{i : \tau_i = \tau,\ \beta_i \in J\},
\]
and $W(S) = \sum_{s \in S} W(s)$, we get:

\begin{lemma}[Hall over intervals]
\label{lem:hall}
A window multiset $W$ admits all $n$ items if and only if
\[
  W\big(\mathrm{ext}(J)\big) \;\ge\; M(J) := \max_{\tau \in [N]} n_\tau(J)
  \qquad\text{for every block interval } J .
\]
\end{lemma}

\begin{proof}[Proof sketch]
Necessity is clear: the items of one bin inside $J$ can only be served by
windows starting in $\mathrm{ext}(J)$. For sufficiency, fix a bin $\tau$ and a
set $S$ of its items; split the neighbourhood $N(S)$ into maximal runs of
consecutive starts. The items of one run lie in a block interval
$J_c = [x_c, y_c]$ whose neighbourhood is exactly $\mathrm{ext}(J_c)$, and the
runs are disjoint, so the supplies add: $|N(S)| \ge \sum_c W(\mathrm{ext}(J_c))
\ge \sum_c n_\tau(J_c) \ge |S|$. Hall's theorem then gives a matching for every
bin, and the matchings for different bins do not interact.
\end{proof}

\subsection{A min--max theorem}

Lemma~\ref{lem:hall} turns sequencing into a covering problem: minimize
$\sum_s W(s)$ subject to $W(\mathrm{ext}(J)) \ge M(J)$ for all intervals $J$,
over integers $W(s) \ge 0$. Each constraint's support $\mathrm{ext}(J)$ is a
set of \emph{consecutive} starts, so the constraint matrix has the
consecutive-ones property and is therefore totally unimodular. Its LP
relaxation has an integral optimum -- so the LP value \emph{is} $L^\star$ --
and the dual is integral as well, which yields:

\begin{theorem}[Min--max for sequencing]
\label{thm:minmax}
\[
  L^\star \;=\; \max\Big\{ \textstyle\sum_{k=1}^{K} M(J_k) \;:\;
     J_1,\dots,J_K \text{ block intervals with } \mathrm{ext}(J_k)
     \text{ pairwise disjoint} \Big\}.
\]
\end{theorem}

The ``$\ge$'' direction is combinatorial and worth stating on its own, since it
is what the security argument leans on: fix such a family; a layer serving an
item of $J_k$'s heaviest bin must start in $\mathrm{ext}(J_k)$, the extensions
are disjoint, so each layer is counted for at most one $k$, and $J_k$ needs
$M(J_k)$ distinct layers. Two special cases are instructive:
\begin{itemize}
  \item $K = 1$ with $J = [1,b]$ gives $L^\star \ge \max_\tau (\text{load of }
        \tau)$, the pigeonhole bound;
  \item chopping $[1,b]$ into $\lceil b/z \rceil$ windows and adding their
        per-window maxima is one feasible \emph{primal}, i.e. an upper bound;
        it is what a naive ``hard partition'' sequencer achieves, and
        Theorem~\ref{thm:minmax} shows it is generally far from optimal.
\end{itemize}

\paragraph{The algorithm.}
We sequence with a single left-to-right first-fit pass: process blocks in
order and place each item in the admissible layer of smallest anchor, opening
a new layer anchored at the current block when none fits. It costs $O(n L)$
-- $14$\,ms at $n = 2^{20}$. (The covering program could also be solved
exactly by the standard interval-covering sweep -- add the missing windows at
the latest admissible start whenever a constraint falls short, then assign
each item to the leftmost admissible window, which is optimal for convex
bipartite matching~\cite{Glover67} -- in $O(nb + b^2)$, $130$\,ms at
$n = 2^{20}$; the point of this subsection is that nothing of the sort is
needed.) Despite its simplicity the first-fit pass is not a heuristic at all:

\begin{theorem}[The first-fit sequencer is exact]
\label{thm:greedy}
On every input the first-fit sequencer opens exactly $L^\star$ layers.
More precisely, writing $G(p)$ for the number of its layers with anchor
$\le p$ and
\[
  D(p) \;=\; \max\Big\{ \textstyle\sum_{k} M(J_k) \;:\;
     J_k = [x_k, y_k] \text{ with } \mathrm{ext}(J_k)
     \text{ pairwise disjoint},\ y_k \le p \Big\},
\]
it achieves $G(p) = D(p)$ for every $p$, whereas \emph{every} feasible
layering $W$ satisfies $\sum_{s \le p} W(s) \ge D(p)$. The first-fit layering
is therefore optimal, and pointwise the latest optimal one.
\end{theorem}

The proof rests on a structural property of the failure events. Throughout,
a layer's \emph{anchor} $a(\ell)$ is the block current when it was opened (so
anchors are created in nondecreasing order, and a layer admits blocks
$[a(\ell), a(\ell)+z-1]$), and ``$\tau$-occupied'' means the layer already
holds an item of bin $\tau$; occupied slots never vacate.

\begin{lemma}[Blocking structure]
\label{lem:blocking}
Suppose the first-fit sequencer, holding the layer multiset $\Lambda$, fails
on an item $t$ of bin $\tau$ at block $y$. Then there is an
$X \le y - z + 1$ such that
(i) the layers with $a(\ell) \ge X$ all have $a(\ell) \in [X, y]$, and each
of them is $\tau$-occupied by an item with block in $[X + z - 1,\, y]$;
(ii) hence, with $K = |\{\ell \in \Lambda : a(\ell) \ge X\}|$ and counting
$t$ itself, the items processed so far satisfy
$M([X + z - 1,\, y]) \ge K + 1$.
\end{lemma}

\begin{proof}
Build a closure. Maintain a collection of \emph{ranges} of anchors, initially
the single range $[y - z + 1,\, y]$; whenever a layer $\ell \in \Lambda$ has
its anchor inside a known range, absorb it: we argue below that $\ell$ is
$\tau$-occupied, by a unique item $w_\ell$ (its \emph{witness}) with block
$\beta_\ell$; add the range $[\beta_\ell - z + 1,\; s_\ell]$, where
$s_\ell = a(\ell)$. Iterate to a fixpoint; let $X$ be the least left endpoint
of any range.

\emph{Every absorbed layer is $\tau$-occupied.} Induct over the absorption
order. A layer absorbed by the initial range was admissible to $t$, and $t$
failed. Otherwise $\ell$'s anchor lies in $[\beta_p - z + 1, s_p]$ for an
earlier witness $p$. If $a(\ell) < s_p$: layers are opened at the block
current at their creation, so $\ell$ was opened while block
$a(\ell) < s_p \le \beta_p$ was current -- before $p$ was processed -- and at
$p$'s turn it was admissible ($a(\ell) \ge \beta_p - z + 1$) and preceded
$p$'s layer in the smallest-anchor-first scan; $p$ passing it over means it
was $\tau$-occupied then, hence still. If $a(\ell) = s_p$: either $\ell$ is
$p$'s own layer, occupied by $p$; or trace how $p$'s layer was absorbed --
through a range containing $s_p$ strictly below its endpoint, in which case
the previous case applies to $\ell$ verbatim; through the initial range, so
$a(\ell) \in [y-z+1, y]$ and $t$'s failure applies; or through another
witness's endpoint $s_q = s_p$, and we ascend to $q$. The ascent follows the
absorption order, so it terminates in one of the resolved cases.

\emph{The ranges cover exactly $[X, y]$.} A new range
$[\beta_\ell - z + 1, s_\ell]$ contains $s_\ell$ (an item's block is within
$z - 1$ of its layer's anchor), which lies in the range that absorbed
$\ell$; so the union stays a single interval, and its right end stays $y$
(no anchor exceeds the current block).

Consequently every layer with $a(\ell) \ge X$ has $a(\ell) \in [X, y]$, is
covered by a range, and was absorbed -- which is (i), since each witness has
$\beta_\ell - z + 1 \ge X$ (its range's left endpoint is one of those $X$
minimizes over), i.e. $\beta_\ell \ge X + z - 1$. For (ii): the witnesses
are distinct (one per layer, one bin-$\tau$ slot per layer), all of bin
$\tau$, with blocks in $[X + z - 1, y]$; together with $t$ (block
$y \ge X + z - 1$) they number $K + 1$.
\end{proof}

\begin{proof}[Proof of Theorem~\ref{thm:greedy}]
First, \emph{localized weak duality}: for any layering $W$ of the items
processed so far and any family as in the definition of $D(p)$, the
$M(J_k)$ items of $J_k$'s heaviest bin occupy $M(J_k)$ distinct windows with
starts in $\mathrm{ext}(J_k) \subseteq (-\infty, y_k] \subseteq
(-\infty, p]$, and the extensions are disjoint, so
$\sum_{s \le p} W(s) \ge \sum_k M(J_k)$. Hence
$\sum_{s \le p} W(s) \ge D(p)$ for every feasible $W$ -- in particular
$G(p) \ge D(p)$ at termination, since the first-fit layering is feasible.

Second, the invariant $G(p) \le D(p)$ for all $p$, by induction over the
processed items ($G$, $D$, $M$ all refer to the current prefix). Placing an
item in an existing layer changes no $G(p)$ and never decreases $D(p)$.
Suppose item $t$ (bin $\tau$, block $y$) triggers a new layer, and let $X$,
$K$ be as in Lemma~\ref{lem:blocking}. For $p < y$ nothing changes: the new
anchor is $y$. For $p \ge y$, take a family with right endpoints
$\le X - 1$ certifying $G(X-1)$ (induction), and adjoin
$J_0 = [X + z - 1,\, y]$: its extension $\mathrm{ext}(J_0) = [X, y]$ is
disjoint from the previous extensions $\subseteq (-\infty, X - 1]$, and
$M(J_0) \ge K + 1$ by the lemma, counting $t$. The family's value is now at
least $G(X - 1) + K + 1$, which by Lemma~\ref{lem:blocking}(i) is exactly
the new $G(p)$.

At termination the two directions give $G(p) = D(p)$ for every $p$; at
$p = b$, the first-fit count $G(b)$ equals $D(b)$, which by weak duality is
at most $L^\star$, while $G(b) \ge L^\star$ by feasibility.
\end{proof}

\begin{remark}
The argument re-proves the min--max theorem constructively: the first-fit
layering is feasible with $D(b)$ windows, so
Theorem~\ref{thm:minmax} follows without invoking total unimodularity.
The same equality $G \equiv D$ shows the first-fit sequencer places windows
at exactly the positions the exact covering sweep would -- they differ in
cost, not output. Both facts are corroborated mechanically: beyond the
$9{,}400$ randomized small instances checked against a brute-force chromatic
number and the production sets checked against the dual, the closure of
Lemma~\ref{lem:blocking} and the profile identity $G \equiv D$ were
machine-verified at every failure event across a further $5 \cdot 10^5$
instances, including exhaustive enumerations of all item placements for
small $(N, b, z, n)$.
\end{remark}

\subsection{The public layer budget}
\label{subsec:budget}

The realized layer count is a function of the sender's set: it is driven by
which items collide in a bin, and collisions are determined by
$\mathsf{H}_1$ evaluated on the sender's items. Transmitting $L^\star$ would
therefore leak. Instead both parties derive a \emph{public budget}
$\bar{L}$ from the parameters alone, and the sender always transmits exactly
$\bar{L}$ layers -- the $L^\star$ real ones, and $\bar{L} - L^\star$ padding
layers whose ciphertexts are a random plaintext multiple of a received
ciphertext, plus the usual mask.

Padding only goes upwards, so the sender aborts exactly when no layering fits
inside the budget, i.e. when $L^\star > \bar{L}$; Theorem~\ref{thm:budget}
bounds that event by $2^{-\lambda}$. The bound is a statement about the
\emph{minimum} layer count, so the sender has to attain it -- a sequencer
that overshot $L^\star$ would abort more often than the theorem allows. By
Theorem~\ref{thm:greedy} the first-fit sequencer attains it on every input,
so the abort event is exactly $\{L^\star > \bar{L}\}$. Conversely, when the budget is met, nothing about the realized count reaches
the wire: the receiver gets $\bar{L}$ ciphertexts and $\bar{L} \cdot N$
equality instances whatever $L^\star$ was, and which algorithm produced the
layering is immaterial.

\begin{remark}[Padding and noise]
Fixing the transmitted count is necessary but not sufficient: a real layer
accumulates $w + z$ ciphertext--plaintext products while a padding layer is a
single one, so their noise magnitudes differ by several bits, and a receiver
holding the decryption key can recover the noise and count the real layers.
Making the layers indistinguishable is the job of the noise flooding applied
before the ciphertexts are returned -- inflating the padding to match, which we
also tried, does not achieve it, since real layers already differ among
themselves by occupancy.
\end{remark}

\begin{theorem}[Layer budget]
\label{thm:budget}
Let the $n$ positions be uniform and independent over $[R]$, $R = m-w+1$, and
let
\[
  D(g) \;=\; \min\Big\{ k \;:\;
     \Pr\big[\,\mathrm{Bin}(n,\, g/R) \ge k \,\big] \;\le\;
     \frac{2^{-\lambda}}{N b^2} \Big\},
  \qquad
  \rho^\star \;=\; \max_{1 \le g \le b} \frac{D(g)+1}{g+z-1}.
\]
Then, with $\bar{L} = \lceil \rho^\star (b+z-1) \rceil$,
\[
  \Pr\big[\, L^\star > \bar{L} \,\big] \;\le\; 2^{-\lambda}.
\]
\end{theorem}

\begin{proof}
Fix a bin $\tau$, a start $x$ and a length $g$. Each item lands in bin $\tau$
with block in $[x, x+g-1]$ iff its position is one of the $g$ positions mapping
there, so $n_\tau([x,x+g-1]) \sim \mathrm{Bin}(n, \le g/R)$ and by definition of
$D(g)$ it exceeds $D(g)$ with probability at most $2^{-\lambda}/(N b^2)$. There
are at most $N b^2$ triples $(\tau, x, g)$, so by a union bound
\begin{equation}
  \label{eq:allJ}
  \Pr\big[\, \exists J:\ M(J) > D(|J|) \,\big] \;\le\; 2^{-\lambda}.
\end{equation}
Assume the complement. Lay windows down at uniform density $\rho^\star$ over
the $b+z-1$ admissible starts $s \in [-z+1, b-1]$, i.e.
$W(s) = \lfloor (s+1)\rho^\star \rfloor - \lfloor s \rho^\star \rfloor$; then
any $\ell$ consecutive starts receive at least $\rho^\star \ell - 1$ windows.
An interval $J$ with $|J| = g$ has $|\mathrm{ext}(J)| = g+z-1$, hence
\[
  W(\mathrm{ext}(J)) \;\ge\; \rho^\star (g+z-1) - 1 \;\ge\; D(g) \;\ge\; M(J),
\]
where the middle inequality is the definition of $\rho^\star$. By
Lemma~\ref{lem:hall} this $W$ is feasible, and it uses
$\lceil \rho^\star (b+z-1) \rceil = \bar{L}$ windows, so
$L^\star \le \bar{L}$.
\end{proof}

Three remarks on how the bound behaves.

\paragraph{It is computed, not extrapolated.}
$D(g)$ is an exact binomial quantile: the upper tail is summed term by term in
log space, so $\bar{L}$ is a closed-form function of $(n, N, b, z, \lambda)$
that both parties evaluate in milliseconds. Nothing here is fitted or
extrapolated from measurements, in contrast with the band-width selection of
the underlying OKVS.

\paragraph{Where the slack is.}
$\bar{L}$ overshoots the realized $L^\star$ for two reasons: the union bound
charges every one of the $Nb^2$ constraints separately, which is expensive for
short intervals where the mean count is small; and the uniform-density primal
pays the worst local ratio everywhere. Both shrink as $z$ grows -- the
denominator $g+z-1$ absorbs the union tax, and the binding constraint moves
towards $J = [1,b]$, where the union is over $N$ bins only. Measured at
$n = 2^{20}$ ($b = 269$):
\[
  \begin{array}{c|cccc}
    z & 20 & 40 & 60 & 80 \\ \hline
    \bar{L} & 399 & 296 & 266 & 255 \\
    \mathbb{E}[L^\star] & 196 & 178.5 & 175.5 & 175.5 \\
    \bar{L} / \mathbb{E}[L^\star] & 2.04 & 1.66 & 1.52 & 1.45
  \end{array}
\]
The pigeonhole floor -- the $2^{-\lambda}$ quantile of the maximum bin load,
which \emph{no} budget can go below -- is $232$ here, so at $z = 80$ the
certificate costs only $10\%$ over the information-theoretic minimum.

\paragraph{Cost of the padding.}
A padding layer costs one plaintext multiplication (versus $w + z$ for a real
layer), so the compute is driven by $L^\star$; the communication and the number
of equality instances are driven by $\bar{L}$. Increasing $z$ therefore trades
a larger per-layer decode against a smaller $\bar{L}$, and in our measurements
the second effect dominates up to $z \approx 80$.
